\documentclass{article}

\usepackage{amsmath,amssymb}
\usepackage{xurl}
\usepackage[colorlinks=true,linkcolor=blue,citecolor=blue,urlcolor=blue]{hyperref}
\setlength{\emergencystretch}{2em}

% Shared commands for notes in docs/paper-gaps/.

\DeclareUnicodeCharacter{2080}{\ifmmode{}_{0}\else\textsubscript{0}\fi}
\DeclareUnicodeCharacter{2081}{\ifmmode{}_{1}\else\textsubscript{1}\fi}
\DeclareUnicodeCharacter{2082}{\ifmmode{}_{2}\else\textsubscript{2}\fi}
\DeclareUnicodeCharacter{2099}{\ifmmode{}_{n}\else\textsubscript{n}\fi}

\newcommand{\Lean}{\textsc{Lean}}
\ExplSyntaxOn
\NewDocumentCommand{\leanid}{m}{
  \ifmmode
    \textsf{\detokenize{#1}}
  \else
    \group_begin:
    \urlstyle{sf}
    \tl_set:Nn \l_tmpa_tl {#1}
    \tl_replace_all:Nnn \l_tmpa_tl {\_} {_}
    \exp_args:NV \path \l_tmpa_tl
    \group_end:
  \fi
}
\ExplSyntaxOff
\providecommand{\tr}{\operatorname{tr}}
\newcommand{\tnleanIssueBase}{https://github.com/LionSR/TNLean/issues/}
\newcommand{\tnleanPrBase}{https://github.com/LionSR/TNLean/pull/}
\newcommand{\tnleanDocBase}{https://sirui-lu.com/TNLean/docs/}
\newcommand{\ghissue}[1]{\href{\tnleanIssueBase#1}{GitHub issue \##1}}
\newcommand{\ghpr}[1]{\href{\tnleanPrBase#1}{PR \##1}}
\newcommand{\doclink}[1]{\href{\tnleanDocBase#1.html}{\path{#1}}}

% Dirac notation and the identity operator, as in blueprint/src/macros/common.tex.
% \providecommand keeps note-local definitions authoritative.
\usepackage{dsfont}
\providecommand{\ket}[1]{|#1\rangle}
\providecommand{\bra}[1]{\langle #1|}
\providecommand{\braket}[2]{\langle #1 | #2 \rangle}
\providecommand{\ketbra}[2]{|#1\rangle\!\langle #2|}
\providecommand{\Id}{\mathds{1}}
\providecommand{\one}{\mathds{1}}
\providecommand{\C}{\mathbb{C}}
\providecommand{\MN}[1]{M_{#1}(\C)}
\providecommand{\MD}{M_D(\C)}

% Verdict marker: \gapnote{<kind>}{<status>}. Kind is the policy
% classification (clarification, local-correction, scope-restriction,
% unfaithful, false-source, open-gap); status is open, wip (elimination
% underway), resolved, or historical. Machine-read by the site index;
% prints a verdict line.
\newcommand{\gapnote}[2]{\par\noindent\textbf{Verdict:} #1 (#2).\par}


% Fallbacks keep this author document compilable when it is built outside its
% source directory and a different command.tex is found first.
\providecommand{\leanid}[1]{\textsf{\detokenize{#1}}}
\providecommand{\doclink}[1]{%
    \href{https://sirui-lu.com/TNLean/docs/#1.html}{\path{#1}}}
\providecommand{\gapnote}[2]{%
    \par\noindent\textbf{Verdict:} #1 (#2).\par}

\title{Safe-Box Entropy and the Sublinear Collar Bound}
\author{}
\date{2026-10-15}

\begin{document}
\maketitle
\gapnote{scope-restriction}{resolved}

\section{The source propositions}

Propositions~9.5 and~9.6 of the two-dimensional area-law manuscript of
September~24, 2026~\cite{OpenAI2026AreaLaw} (\path{08-scanner.tex},
lines~681--944, at revision \texttt{adc7f1241b42e322a6451854ab7e4b4c146bf78a})
establish the small-exponent entropy bound on safe rectangles and deduce from
it the sublinear mutual-information bound across the collar of a separated
template.

Fix a local Hilbert-space dimension $q \ge 1$, a finite interaction range
$R \in \mathbb N$, an interaction norm bound $J \ge 0$, and a spectral gap
$\Delta > 0$. Consider a finite lattice domain $\Lambda \subset \mathbb Z^2$, a
local Hamiltonian on $\Lambda$ with those parameters and ground energy $e_0$,
a unit ground vector $\Omega$, and a subset $A \subseteq \Lambda$ with cut
boundary $\partial A$ in the induced graph on $\Lambda$. Write $Z$ for
the set of endpoints of these cut edges, and interpret distance to the cut
as distance to $Z$, with distance to the empty set equal to infinity. For any region
$Y \subseteq \Lambda$, write $S_\Omega(Y)$ for the von Neumann entropy of the
reduced density matrix of $\Omega$ on $Y$.

Fix a safety parameter $D_0 > 2R + 10$. A lattice rectangle
$Q \subset \mathbb Z^2$ with nonempty coordinate intervals has size
$s = \operatorname{size}(Q)$ equal to the larger of its two side lengths in
lattice sites. The rectangle $Q$ is $D_0$-safe for the cut $A$ if the
$\ell_\infty$ distance $\operatorname{dist}_\infty(Q, \partial A)$ exceeds
$D_0 s$. For each integer scale $r \ge 1$, let $F_{\mathrm{box}}(r)$ denote the
supremum of $S_\Omega(A \cap Q)$ over all finite domains $\Lambda$, gapped
ground vectors $\Omega$, cuts $A \subseteq \Lambda$, and $D_0$-safe rectangles
$Q$ of size at most $r$. Proposition~9.5 of Ref.~\cite{OpenAI2026AreaLaw}
(\path{08-scanner.tex}, lines~681--805) states that for the target exponent
$e_* = 2 \times 10^{-6}$, there exists a constant $C_{\mathrm{box}} > 0$,
depending only on $q$, $R$, $J$, and $\Delta$, such that for every
$D_0 > 2R + 10$ and every integer $r \ge 1$,
\begin{align}
    F_{\mathrm{box}}(r)
     & \le C_{\mathrm{box}} r^{1 + e_*}.
    \label{eq:safe_box_collar_prop95}
\end{align}
The paper proves (\ref{eq:safe_box_collar_prop95}) from an initial safe-box
exponent $0 < e_0 < 1$ by iterating the one-step exponent map
$e \mapsto (1 - \ell_0)e$, where $\ell_0 = (1 - e_0)/200000$, finitely many
times until the exponent is at most $e_*$.

Proposition~9.6 of Ref.~\cite{OpenAI2026AreaLaw} (\path{08-scanner.tex},
lines~808--944) applies (\ref{eq:safe_box_collar_prop95}) to a template
$T \subset \mathbb Z^2$ at scanner scale $n \ge 2$ and template scale $s_0$.
With the collar exponent $\ell = 10^{-5}$ and collar depth
$L = \lfloor n^{1 - \ell}\rfloor$, let $T_L$ denote the $\ell_\infty$
neighborhood of $T$ of radius $L$, and define the core region $X = A \cap T$
and the collar truncation set $S_0 = A \cap T_L$. For a template $T$ satisfying
$\operatorname{dist}_\infty(T, \partial A) > 4 D_0 s_0$ and $L \le s_0$,
Proposition~9.6 bounds the mutual information $I_\Omega(X : \Lambda \setminus S_0)$
between $X$ and $\Lambda \setminus S_0$ by
\begin{align}
    2 I_\Omega(X : \Lambda \setminus S_0)
     & \le C n^{1 - 1/100000},
    \label{eq:safe_box_collar_prop96}
\end{align}
for a constant $C > 0$ and scale threshold $N \ge 2$ chosen after the fixed
system and template parameters and before the finite domain $\Lambda$,
Hamiltonian, ground vector $\Omega$, cut $A$, and template $T$.

\section{The formal statement and hypothesis reduction}

The conditional template mutual-information theorem derives
(\ref{eq:safe_box_collar_prop96}) from the analytic estimate
(\ref{eq:safe_box_collar_prop95}). The latter is now proved independently,
so it is not an additional hypothesis of the physical collar theorem.
For the conditional argument, fix $q \ge 1$, $R \in \mathbb N$,
$J \ge 0$, $\Delta > 0$, $D_0\in\mathbb N$ with $D_0 \ge 1$,
and a template constant
$C_{\mathrm{tpl}} \ge 24$. Assuming (\ref{eq:safe_box_collar_prop95}) at
$e_* = 2 \times 10^{-6}$, there exist $C > 0$ and $N \ge 2$ such that every
finite instance with $n \ge N$,
$\operatorname{dist}_\infty(T, \partial A) > 4 D_0 s_0$, and
$L = \lfloor n^{1 - 1/100000}\rfloor \le s_0$ satisfies
\begin{align}
    2 I_\Omega(X : \Lambda \setminus S_0)
     & \le C n^{62497/62500}
    \le C n^{1 - 1/100000}.
    \label{eq:safe_box_collar_conditional_template}
\end{align}
All intermediate hypotheses of the scanner chain are eliminated in five steps
leading to (\ref{eq:safe_box_collar_conditional_template}).\footnote{The
    conditional template mutual-information theorem is
    \leanid{TNLean.PEPS.AreaLaw.Scan.exists_template_final_mutualInformation_of_boxEntropy},
    tracked in \href{https://github.com/LionSR/TNLean/issues/8756}{issue~8756}
    and \href{https://github.com/LionSR/TNLean/issues/9334}{issue~9334}.}

\paragraph{1. Post-replica collar mutual information from the sharp comparison.}
Fix the final scanner exponents $\mu = 1/4$, $\ell = 10^{-5}$,
$\omega = 1/20000$, $\kappa = 10^{-6}$, and $\nu = 1/1000$, and define the
integer scales
\begin{align}
    L   & = \lfloor n^{1 - \ell}\rfloor,             &
    m   & = \lfloor n^\mu\rfloor,                    &
    D   & = \lceil n^\kappa\rceil, \nonumber           \\
    K   & = L \mathbin{\mathrm{div}} (8m),           &
    r_0 & = \lceil C_{\mathrm{tr}} (\log n)^2\rceil, &
    a   & = \frac{W}{K},
    \label{eq:safe_box_collar_scales}
\end{align}
where $m$ is the band spacing, $D$ is the separation scale, $K$ is the band
count, $r_0$ is the truncation radius with coefficient $C_{\mathrm{tr}} > 0$,
$W \in [1, n^\omega]$ is the total band weight, and $a$ is the per-band weight
(\path{08-scanner.tex}, lines~32--58 and~839--865). Let $\Omega_t$ be the unit
ground vector of the Hamiltonian truncated to $S_0$, with ground energy
$e_{\mathrm{tr}} \le \min(n^{-1000}, g/4)$, normalized gap $g/2$ where
$g = \Delta / \max\{1, J(\|f_{1,\Delta/2}\|_1 + 1)\}$, and trace distance
$d_{\mathrm{tr}}(\Omega_t, \Omega) \le \sqrt{2\min(n^{-1000}, g/4)/g}$. Let
$E$ be a typical spectral window of the reduced state of $\Omega_t$ on $X$ with
dimension $d = |E|$, spectral mass $z > 0$, and auxiliary Schmidt entropy $S'$
satisfying $|S' - S_{\Omega_t}(X)| \le c$ for $c = n^{3/5} - \log z$.

Assume the target and positive-depth prefix entropy bounds on the original
ground vector $\Omega$ at exponent $e_{\mathrm{box}} = e_*$:
\begin{align}
    S_\Omega(X)
     & \le C_e n^{1 + e_*},               &
    S_\Omega(Q_j)
     & \le C_e \bigl(n L^{e_*} + n\bigr),
    \label{eq:safe_box_collar_entropy_inputs}
\end{align}
for every positive-depth prefix $Q_j$ and a constant $C_e \ge 0$. Regional
continuity transfers (\ref{eq:safe_box_collar_entropy_inputs}) to $\Omega_t$
with shell bound $B_{\mathrm{sh}} = C_e(n L^{e_*} + n) + n^{-247}$. At replica
count $k$, averaging over the $n m$ charge rounds and the parameter interval
$[n^{-\nu}, 2n^{-\nu}]$ selects a charge round $s_k$ and parameter
$p_k \in [n^{-\nu}, 2n^{-\nu}]$ (a rescaling by two of the paper's interval
$[n^{-\nu}/2, n^{-\nu}]$ in \path{08-scanner.tex}, lines~380--384 and~880--883).
Because $e_* - \mu + \nu < -8\nu$ (with margin $120499/500000$) and
$(\omega - (1 - \ell - \mu))/4 + \nu < -8\nu$ (with margin $35697/200000$), the
normalized charge entropy defect $\mathcal Q_k(s_k, p_k)/(K n D)$ at the
selected point is at most $n^{-8\nu} = n^{-1/125}$ for large $k$. (An
intermediate formal lemma recorded the weaker bound $n^{-\nu}$, which is too
weak after taking the eighth root in the energy estimate; the paper's bound
$\delta_n \le n^{-0.18}$ in \path{08-scanner.tex}, lines~880--883, is stronger
than $n^{-8\nu}$ and needs no modification.) Substituting
$(\mathcal Q_k/(K n D))^{1/8} \le n^{-\nu}$ and $p_k \le 2n^{-\nu}$ into the
selected-point energy cap with exponent
$2\omega + \ell + 2\kappa - \nu = -111/125000 < -1/2000$ gives the normalized
energy defect bound
\begin{align}
    0 \le \delta_k
     & \le n^{-1/2000} + \frac{\mathrm{rem}_k}{g/2},
    \label{eq:safe_box_collar_energy_defect}
\end{align}
where $\mathrm{rem}_k \to 0$ as $k \to \infty$ (\path{08-scanner.tex},
lines~457--533 and~839--911).\footnote{The moment-margin declarations are
    \leanid{TNLean.PEPS.AreaLaw.Scan.CollarScan.exists_truncated_target_actualSharpUpper}
    in \doclink{TNLean/PEPS/AreaLaw/Scan/SourceTargetSharpUpper} and
    \leanid{TNLean.PEPS.AreaLaw.Scan.CollarScan.exists_truncated_target_sharp_moments}
    in \doclink{TNLean/PEPS/AreaLaw/Scan/SourceTargetSharpMoments}; the
    endpoint budget and entropy-transfer declarations carrying
    (\ref{eq:safe_box_collar_entropy_inputs}) are
    \leanid{TNLean.PEPS.AreaLaw.Scan.CollarScan.eventually_actualScan_suppliedPrevector_endpoint_budget}
    in \doclink{TNLean/PEPS/AreaLaw/Scan/ActualEndpointRoughBudget},
    \leanid{TNLean.PEPS.AreaLaw.Scan.CollarScan.exists_source_actualScan_suppliedPrevector_endpoint_budget}
    in \doclink{TNLean/PEPS/AreaLaw/Scan/SourceEndpointRoughBudget},
    \leanid{TNLean.PEPS.AreaLaw.Scan.eventually_abs_regionalEntropy_sub_le_of_truncation}
    in \doclink{TNLean/PEPS/AreaLaw/Scan/TargetEntropyTransfer},
    \leanid{TNLean.PEPS.AreaLaw.Scan.roughBudget_le_source_power} and
    \leanid{TNLean.PEPS.AreaLaw.Scan.selectedDensity_budget_le_principal}
    in \doclink{TNLean/PEPS/AreaLaw/Scan/SourceRoughBudgetPower}, and
    \leanid{TNLean.PEPS.AreaLaw.Scan.exists_final_selectedDensity_le}
    in \doclink{TNLean/PEPS/AreaLaw/Scan/SelectedDensityPower}; the density and
    energy exponent-margin declarations are
    \leanid{TNLean.PEPS.AreaLaw.Scan.ScannerExponents.eventually_principalDensity_le_rpow},
    \leanid{TNLean.PEPS.AreaLaw.Scan.ScannerExponents.exists_principalDensity_le_rpow},
    \leanid{TNLean.PEPS.AreaLaw.Scan.final_principalDensity_margins},
    \leanid{TNLean.PEPS.AreaLaw.Scan.exists_final_principalDensity_le_rpow},
    \leanid{TNLean.PEPS.AreaLaw.Scan.final_eightNuDensity_margins},
    \leanid{TNLean.PEPS.AreaLaw.Scan.exists_final_principalDensity_le_eightNu},
    and \leanid{TNLean.PEPS.AreaLaw.Scan.exists_final_selectedDensity_le_eightNu}
    in \doclink{TNLean/PEPS/AreaLaw/Scan/StrongSelectedDensityPower}, together
    with \leanid{TNLean.PEPS.AreaLaw.Scan.finalSelectedEnergyCap_le},
    \leanid{TNLean.PEPS.AreaLaw.Scan.final_selectedEnergy_margins},
    \leanid{TNLean.PEPS.AreaLaw.Scan.eventually_final_selectedEnergy_le_rpow},
    \leanid{TNLean.PEPS.AreaLaw.Scan.exists_final_selectedEnergy_le_rpow}, and
    \leanid{TNLean.PEPS.AreaLaw.Scan.exists_final_selectedEnergy_add_remainder_le_rpow};
    the composed selected-point comparison, density, and energy bounds are
    \leanid{TNLean.PEPS.AreaLaw.Scan.CollarScan.exists_sourceTarget_selected_actualCharge_sharp_comparison},
    \leanid{TNLean.PEPS.AreaLaw.Scan.CollarScan.exists_sourceTarget_selected_actualCharge_sharp_density}
    in \doclink{TNLean/PEPS/AreaLaw/Scan/SourceSelectedSharpDensity}, and
    \leanid{TNLean.PEPS.AreaLaw.Scan.CollarScan.exists_sourceTarget_selected_actualCharge_sharp_energy}
    in \doclink{TNLean/PEPS/AreaLaw/Scan/SourceSelectedSharpEnergy}.}

Subtracting the sharp upper and lower norm rates at $(s_k, p_k)$, substituting
(\ref{eq:safe_box_collar_energy_defect}), bounding the terminal pre-charge
entropies on $\Omega_t$ by $2 S_{\Omega_t}(X) + 2 B_{\mathrm{sh}} + 6 n D \log q$,
transferring mutual information from $\Omega_t$ to $\Omega$ with the continuity
term $C_{\mathrm{MI}} = 16 \sqrt{d_{\mathrm{tr}}(\Omega_t, \Omega)} \log(e \dim \mathcal H_X)$,
and taking $k \to \infty$ at fixed $n$, $W \in [1, n^\omega]$, and smoothing
parameter $\tau \in (0, 1/2]$ yields
\begin{align}
    2 I_\Omega(X : \Lambda \setminus S_0)
     & \le \frac{4 S_{\Omega_t}(X) + 4c - 2\log z}{W} \nonumber                     \\
     & \quad + \left(\tau + \frac{n^{-1/2000}}{\tau}\right)
    \bigl(4 S_{\Omega_t}(X) + 2c + 2 B_{\mathrm{sh}} + 6 n D \log q\bigr) \nonumber \\
     & \quad + \left(3 + \frac{1}{a}\right) h(\tau)
    + 37 C_0 B a - \frac{\log z}{a} + 2c + C_{\mathrm{MI}},
    \label{eq:safe_box_collar_post_replica}
\end{align}
where $h(\tau)$ is the binary entropy and $C_0 B$ is the quadratic marginal
moment coefficient with $B = C_B n D (\log n)^{12}$ (\path{08-scanner.tex},
lines~868--911 and~936--941).\footnote{The mutual-information comparison and
    replica limit declarations are
    \leanid{TNLean.PEPS.AreaLaw.Scan.two_information_le_sharpInformationBound}
    and
    \leanid{TNLean.PEPS.AreaLaw.Scan.information_le_of_eventually_sharpInformationBound}
    in \doclink{TNLean/PEPS/AreaLaw/Scan/SharpInformationScalar}, and
    \leanid{TNLean.PEPS.AreaLaw.Scan.CollarScan.exists_sourceTarget_actual_mutualInformation_bound}
    in \doclink{TNLean/PEPS/AreaLaw/Scan/SourceActualInformation}.}

\paragraph{2. Final-scale error absorption and transport smallness.}
Specializing (\ref{eq:safe_box_collar_post_replica}) to $W = n^{1/20000}$ and
$\tau = n^{-1/4000}$, the transport smallness condition
$a \ell_q(r_0) \le c_t$, where $\ell_q(r_0) = 1 + (2r_0 + 1)^2 \log q$, holds
for all sufficiently large $n$ because $\ell_q(r_0) = O((\log n)^4)$ and
$a \le C_a n^{-0.74994}$. From $X \subseteq S_0$ and $|S_0| \le (C_T + 1)n^2$,
the target dimension factor obeys
$\log(e \dim \mathcal H_X) \le 1 + ((C_T + 1)\log q)n^2$, so
$C_{\mathrm{MI}} = O(n^{-248})$. Under (\ref{eq:safe_box_collar_entropy_inputs}),
the seven error terms in (\ref{eq:safe_box_collar_post_replica}) have exponents
\begin{align}
     & 0.999952,\quad 0.999752,\quad 0.74969,\quad 0.250061, \nonumber \\
     & {-99.25006},\quad 0.6,\quad {-248},
    \label{eq:safe_box_collar_error_exponents}
\end{align}
whose maximum in (\ref{eq:safe_box_collar_error_exponents}) is
$62497/62500 = 0.999952 = 1 - 1/100000 - 19/500000$
(\path{08-scanner.tex}, lines~839--942). Absorbing the logarithmic factors into
that margin yields (\ref{eq:safe_box_collar_conditional_template}) conditionally
on (\ref{eq:safe_box_collar_entropy_inputs}), the scale inequalities
$2r_0 \le D$, $1 \le D$, and $4D \le m$, and the scanner geometry.\footnote{The
    final-scale scalar absorption and transport smallness declarations are
    \leanid{TNLean.PEPS.AreaLaw.Scan.finalSharpError},
    \leanid{TNLean.PEPS.AreaLaw.Scan.exists_finalSharp_source_error_le}, and
    \leanid{TNLean.PEPS.AreaLaw.Scan.exists_finalSharp_transport_smallness};
    the target dimension bound is
    \leanid{TNLean.PEPS.AreaLaw.Scan.log_exp_configuration_card_le_of_card}
    in \doclink{TNLean/PEPS/AreaLaw/Scan/PhysicalInformationDimension}; the
    conditional final-scale mutual-information theorem is
    \leanid{TNLean.PEPS.AreaLaw.Scan.CollarScan.exists_sourceTarget_finalScale_mutualInformation}
    in \doclink{TNLean/PEPS/AreaLaw/Scan/SourceFinalInformation}.}

\paragraph{3. Numerical admissibility of the rounded scales.}
At the rounded definitions (\ref{eq:safe_box_collar_scales}), the inequalities
$2r_0 \le D$, $1 \le D$, and $4D \le m$ hold for all sufficiently large $n$
because $r_0 = O((\log n)^2)$, $D = \lceil n^{10^{-6}}\rceil$, and
$m = \lfloor n^{1/4}\rfloor$. This removes the three numerical scale premises,
leaving only the scanner geometry and the entropy bounds
(\ref{eq:safe_box_collar_entropy_inputs}).\footnote{The rounded-scale
    mutual-information theorem is
    \leanid{TNLean.PEPS.AreaLaw.Scan.CollarScan.exists_sourceTarget_finalScale_mutualInformation_of_rounded_scales}.}

\paragraph{4. Target and prefix entropies from the safe-box bound.}
For any separated template $T$ with
$\operatorname{dist}_\infty(T, \partial A) > 4 D_0 s_0$ and $L \le s_0 \le n$,
the dyadic safe-box partitions of $T$ and of each positive-depth shell cover
$X = A \cap T$ and every prefix $Q_j$ by $D_0$-safe rectangles
(\path{08-scanner.tex}, lines~574--590, 641--667, and~828--844). Subadditivity
of entropy then derives both bounds in (\ref{eq:safe_box_collar_entropy_inputs})
from the single safe-box bound (\ref{eq:safe_box_collar_prop95}) at
$e_* = 2 \times 10^{-6}$.\footnote{The template entropy reductions are
    \leanid{TNLean.PEPS.AreaLaw.Scan.exists_template_source_entropy_bounds}
    and
    \leanid{TNLean.PEPS.AreaLaw.Scan.exists_template_final_source_entropy_bounds_of_boxEntropy}
    in \doclink{TNLean/PEPS/AreaLaw/Scan/TemplateSourceEntropy}.}

\paragraph{5. Template scanner geometry.}
Constructing the scanner on the separated template $T$ with $C_{\mathrm{tpl}} \ge 24$
and $D_0 \ge 1$ supplies the induced domain graph, support anchors, quadratic
volume bound $|S_0| \le (C_T + 1)n^2$, layer cardinalities, anchor
multiplicity, and boundary clearance, reducing all hypotheses of
(\ref{eq:safe_box_collar_conditional_template}) to
(\ref{eq:safe_box_collar_prop95}).

\section{The physical safe-box bootstrap}

The physical safe-box estimate is proved first for integer clearance.
For fixed $q,R,J,\Delta$, a single $C_b>0$ works for every integer
$d>2R+10$ and every positive integer scale.\footnote{The theorem is
    \leanid{TNLean.PEPS.AreaLaw.Scan.exists_boxEntropy_le_boxError}.}
Its composition with the conditional template theorem gives the original-state
information bound at integer clearance, without assuming an entropy estimate.

In the bootstrap argument (\path{08-scanner.tex}, lines~681--805), given an
initial safe-box exponent $0 < e_0 < 1$, the bootstrap fixes
$g_0 = (1 - e_0)/2$, $\mu = 1 - g_0$, $\nu = g_0/1000$, $\ell = \nu/100$,
$\kappa = \min\{\ell, e_*/4\}$, and total band weight $W = 4$, with $a = 4/K$
and $K$ of order $n^{g_0 - \ell}$. The sharp comparator cannot be used at these
scales because its marginal-moment condition $a \sqrt{B} \le c_0$ with $B$ of
order $n^{1+\kappa}$ would require a negative exponent in
\begin{align}
    \frac{1 + \kappa}{2} - g_0 + \ell
     & = \frac{e_0}{2} + \frac{\kappa}{2} + \ell > 0.
    \label{eq:safe_box_collar_bootstrap_obstruction}
\end{align}
Because the exponent in (\ref{eq:safe_box_collar_bootstrap_obstruction}) is
strictly positive, the bootstrap instead uses the rough lower norm comparison at
$W = 4$, $K \ge 8$, and $a = 4/K$, which requires no marginal-moment bound and
gives, for every $0 < \tau \le 1/2$,
\begin{align}
    4S' - 6\left(\tau \log d + \frac{h(\tau)}{a}\right)
    - 3\frac{\delta_k}{\tau}\max\{0, 2S'\} - R_k
     & \le -\frac{1}{k}\log \mathcal N_k,
    \label{eq:safe_box_collar_rough_lower}
\end{align}
where $\mathcal N_k$ is the squared filtered norm and $R_k \to 0$ as
$k \to \infty$ (\path{07-comparators.tex}, lines~585--651;
\path{08-scanner.tex}, lines~755--769). On each $D_0$-safe rectangle $Q$ of
size $s$, the scanner is specialized to $T = Q \cap \mathbb Z^2$ and $n = 24s$,
where the weak clearance $\operatorname{dist}_\infty(Q, \partial A) > D_0 s$
implies $2L + 10r_0 \le D_0 s$ and $(D_0 + 1)L \le D_0 s$ for large $s$
(\path{08-scanner.tex}, lines~693--755).\footnote{The bootstrap scale numerics,
    weak-rectangle scanner geometry, and rough lower comparison are
    \leanid{TNLean.PEPS.AreaLaw.Scan.BootstrapParameters.scannerExponents}
    in \doclink{TNLean/PEPS/AreaLaw/Scan/BootstrapScannerExponents},
    \leanid{TNLean.PEPS.AreaLaw.Scan.ScannerExponents.exists_fixedWeightFour_transport_smallness}
    in \doclink{TNLean/PEPS/AreaLaw/Scan/FixedWeightTransportSmallness},
    \leanid{TNLean.PEPS.AreaLaw.Scan.ScannerExponents.exists_rounded_admissibility}
    in \doclink{TNLean/PEPS/AreaLaw/Scan/ScannerRoundedAdmissibility},
    \leanid{TNLean.PEPS.AreaLaw.Scan.ScannerExponents.exists_rounded_clearance_threshold}
    in \doclink{TNLean/PEPS/AreaLaw/Scan/RoundedClearance},
    \leanid{TNLean.PEPS.AreaLaw.Scan.rectangleBootstrapScanner} and
    \leanid{TNLean.PEPS.AreaLaw.Scan.exists_rectangleBootstrapScanner_geometry}
    in \doclink{TNLean/PEPS/AreaLaw/Scan/RectangleBootstrapGeometry},
    \leanid{TNLean.PEPS.AreaLaw.Scan.CollarScan.exists_actualScan_rough_lower_norm}
    in \doclink{TNLean/PEPS/AreaLaw/Scan/ActualRoughLower},
    \leanid{TNLean.PEPS.AreaLaw.Scan.CollarScan.exists_actualScan_suppliedPrevector_rough_lower}
    in \doclink{TNLean/PEPS/AreaLaw/Scan/ActualRoughLowerPrevector}, and
    \leanid{TNLean.PEPS.AreaLaw.Scan.tendsto_roughLowerRemainder_of_label_dimension}
    in \doclink{TNLean/PEPS/AreaLaw/Scan/ActualRoughLowerRemainder}.} Combining
(\ref{eq:safe_box_collar_rough_lower}) with the rough upper bound at $W = 4$
establishes the one-step improvement $e \mapsto (1 - \ell_0)e$ for
$e_* < e \le e_0$, and finite iteration yields the integer-clearance safe-box estimate.

\section{Arbitrary real clearance}

The source allows every real $D_0>2R+10$, whereas the safe-box bootstrap
uses integer clearance. Put $d=2R+10$. The integer estimate at $d+1$
can be transferred to $d$ without changing the domain, Hamiltonian, or
ground vector. Indeed, partition a $d$-safe rectangle of size $s\geq2$
into at most four rectangles of size at most $\lceil s/2\rceil$.
Since $d\geq10$, one has
$(d+1)\lceil s/2\rceil\leq ds$, so each child is $(d+1)$-safe.
Subadditivity and the volume bound for the remaining small sizes give
\begin{align}
    F_{\mathrm{box},d}(r)
     & \leq C_w r^{1+e_*},
       & C_w & =4C_b+4\log q+1.
    \label{eq:safe_box_collar_clearance}
\end{align}
Here $C_b$ is the physical coefficient at clearance $d+1$; neither it nor
$C_w$ depends on the subsequently chosen real $D_0$.

Every rectangle satisfying the source's strict distance inequality at
$D_0>d$ is $d$-safe. Thus (\ref{eq:safe_box_collar_clearance}) gives
$S_\Omega(A\cap Q)\leq C_w\operatorname{size}(Q)^{1+e_*}$ for every
such rectangle. This pointwise statement is equivalent to
(\ref{eq:safe_box_collar_prop95}): for rectangles of size at most $r$,
monotonicity of the positive power bounds every entropy by $C_wr^{1+e_*}$;
conversely, take $r=\operatorname{size}(Q)\geq1$. The formal theorem
states the pointwise bound and introduces neither a second notion of safety
nor another entropy supremum.\footnote{The clearance comparison is
    \leanid{TNLean.PEPS.AreaLaw.boxEntropy_power_le_of_succ_clearance};
    the real-clearance conclusion is
    \leanid{TNLean.PEPS.AreaLaw.exists_regionalEntropy_safe_rect_le_boxError}.}

For Proposition~9.6, apply the conditional template information theorem
at the integer clearance $d$ with (\ref{eq:safe_box_collar_clearance}).
It supplies $C_m>0$ and a common threshold before every finite physical
instance. A template separated at $4D_0s_0$ is also separated at $4ds_0$.
Consequently the information bound applies to the same original vector
and the unchanged regions $X=A\cap T$ and $S_0=A\cap T_L$.
Dividing its coefficient by two gives
$I_\Omega(X:\Lambda\setminus S_0)\leq(C_m/2)n^{1-1/100000}$.
The cut and volume bounds then supply all clauses of the sublinear
collar bound with the single coefficient $\max(C_m/2,10)$.\footnote{The
    theorem is \leanid{TNLean.PEPS.AreaLaw.sublinearCollarBound_real}.}
This holds for every real $D_0>2R+10$ and $C_{\mathrm{tpl}}\geq24$,
with no assumed safe-box estimate or mutual-information inequality.

This resolves the safe-box and collar restrictions described here.
The two-family repair construction of Proposition~11.2 and its use in
Theorem~1.1 are proved separately.\footnote{The theorems are
    \leanid{TNLean.PEPS.AreaLaw.twoFamilyDecomposition} and
    \leanid{TNLean.PEPS.AreaLaw.uniformAreaLaw}.}

\bibliographystyle{alpha}
\bibliography{references}

\end{document}
